\documentclass[border=5pt]{standalone}
\usepackage{circuitikz}
\begin{document}
\begin{circuitikz}

  % --- node coordinates ---

  % D1  (diode)
  \draw (5.6000,-9.6000) to[diode, l=$D1$] (5.6000,-12.8000);

  % Q1  (pnp)
  \node [pnp, rotate=270, xscale=-1] at (3.2000,-8.0000) (Q1) {};
  \coordinate (Q1_C) at (8.0000,-4.8000);
  \coordinate (n0) at (8.0000,-4.8000);
  \coordinate (Q1_B) at (5.6000,-8.0000);
  \coordinate (n1) at (5.6000,-8.0000);
  \coordinate (Q1_E) at (3.2000,-4.8000);
  \coordinate (n2) at (3.2000,-4.8000);

  % V1  (voltage)
  \draw (11.2000,-8.0000) to[vsource, l=$V1$] (11.2000,-12.0000);

  % --- connections ---
  \draw (5.6000,-9.6000) -- (5.6000,-8.0000);
  \draw (5.6000,-12.8000) -- (5.6000,-15.2000);
  \draw (8.0000,-4.8000) -- (11.2000,-4.8000);
  \draw (5.6000,-8.0000) -- (5.6000,-9.6000);
  \draw (3.2000,-4.8000) -- (5.6000,-4.8000) -- (5.6000,-12.8000);
  \draw (11.2000,-8.0000) -- (11.2000,-4.8000);
  \draw (11.2000,-12.0000) -- (11.2000,-15.2000);

  % --- wire segments ---
  \draw (5.6000,-15.2000) -- (3.2000,-15.2000);
  \draw (11.2000,-15.2000) -- (5.6000,-15.2000);

  % --- grounds ---

\end{circuitikz}
\end{document}
